% Compile with: pdflatex PTS_rk_interprets_Tarski_Mostowski_Robinson_arithmetic_R_harmonized.tex
% Run twice to resolve references.
\documentclass[11pt,a4paper]{article}
\usepackage[T1]{fontenc}
\usepackage{lmodern,amsmath,amssymb,amsthm,mathtools,microtype,booktabs}
\usepackage[margin=25mm]{geometry}
\usepackage{needspace}
\usepackage[hidelinks]{hyperref}
\newcommand{\rk}[1]{\mathsf{rk}_{#1}}
\newcommand{\PTS}{\mathsf{PTS}_{rk}}
\newcommand{\Rt}{\widetilde{\mathsf R}}
\newcommand{\R}{\mathsf R}
\newcommand{\Good}{\mathsf{Good}}
\newcommand{\Case}{\mathsf{Case}}
\newcommand{\Inv}{\mathsf{Inv}}
\newcommand{\ev}{\operatorname{ev}}
\newcommand{\Add}{\mathsf{Add}}
\newcommand{\Mul}{\mathsf{Mul}}
\newcommand{\z}{\mathsf z}
\newcommand{\suc}{\mathsf s}
\newcommand{\nc}{\mathop{\neg}\nolimits_c}
\newcommand{\ac}{\mathbin{\wedge_c}}
\newcommand{\oc}{\mathbin{\vee_c}}
\newcommand{\iffc}{\mathbin{\leftrightarrow_c}}
\newcommand{\Up}{\mathsf{Up}}
\newcommand{\Da}{D_{\!a}}
\newcommand{\za}{\mathsf z_a}
\newcommand{\sa}{\mathsf s_a}
\newcommand{\adda}{\mathsf a_a}
\newcommand{\ma}{\mathsf m_a}
\newcommand{\Ga}{G_a}
\newcommand{\Ea}{E_a}
\newcommand{\Oa}{O_a}
\newcommand{\Gooda}{\mathsf{Good}_a}
\newcommand{\Casea}{\mathsf{Case}^{a}}
\newcommand{\Inva}{\mathsf{Inv}^{a}}
\newcommand{\nua}[1]{\nu^{a}_{#1}}
\newtheorem{lemma}{Lemma}
\newtheorem{theorem}[lemma]{Theorem}
\newtheorem{proposition}[lemma]{Proposition}
\theoremstyle{remark}\newtheorem*{remark}{Remark}
\setlength{\parskip}{2pt}
\emergencystretch=1.5em
\title{\bfseries Rank-indexed pure type system $\PTS$\newline
interprets Tarski--Mostowski--Robinson\newline arithmetic $\R$}
\author{}\date{13 October 2026}
\hypersetup{pdftitle={Rank-indexed pure type system PTS-rk interprets Tarski-Mostowski-Robinson arithmetic R},
 pdfsubject={Classical relational and functional arithmetic interpretations using dependent function types and beta-conversion}}
\begin{document}
\begin{center}
{\Large\bfseries Rank-indexed pure type system $\PTS$\\[3pt]
interprets Tarski--Mostowski--Robinson arithmetic $\R$\par}
\vspace{8pt}
13 October 2026
\end{center}
\vspace{-0.3em}
\begin{abstract}
We construct a fixed interpretation of Tarski--Mostowski--Robinson
arithmetic $\R$ in $\PTS$. A unary Church domain and certified finite observations
first interpret its relational form $\Rt$. Definable graphs then yield
$\R$; a direct interpretation with arithmetic operators is also given.
All constructions use only dependent function types and $\beta$-conversion.
\end{abstract}
\setlength{\abovedisplayskip}{8pt plus 2pt minus 2pt}
\setlength{\belowdisplayskip}{8pt plus 2pt minus 2pt}

\section{Classical arithmetic and the target calculus}\label{sec:logic}
We work in classical first-order logic with equality. Let $\Rt$ be the
relational presentation of Tarski--Mostowski--Robinson arithmetic in the
language $\{0,S,=,\le,A,M\}$, with ternary predicates $A,M$ and numerals
$\bar n=S^n0$. Its axioms are the universal closures of the following schemes,
indexed by external $m,n\in\mathbb N$:
\begin{align}
 A(\bar m,\bar n,x)&\leftrightarrow \overline{m+n}=x,
    \tag{$\widetilde R1$}\label{rtilde:add}\\
 M(\bar m,\bar n,x)&\leftrightarrow \overline{mn}=x,
    \tag{$\widetilde R2$}\label{rtilde:mul}\\
 \bar m&\ne\bar n\quad(m\ne n),
    \tag{$\widetilde R3$}\label{rtilde:neq}\\
 x\le\bar n&\leftrightarrow\bigvee_{i=0}^{n}x=\bar i,
    \tag{$\widetilde R4$}\label{rtilde:bound}\\
 x\le\bar n&\lor\bar n\le x.
    \tag{$\widetilde R5$}\label{rtilde:compare}
\end{align}
Here $m+n$ and $mn$ are calculated externally: the language has no binary
arithmetic function symbols. Section~\ref{sec:nonrelational} relates this theory
to the original, functional arithmetic $\R$~\cite{tmr};
Sections~\ref{sec:operators}--\ref{sec:functional} give a direct interpretation
of those functions.

The pure type system $\PTS$ has sorts $\rk i$ for $i\in\mathbb N$, axioms
$\rk i:\rk0$, and the formation rule
\[
 \frac{\Gamma\vdash T:\rk a\qquad\Gamma,x:T\vdash U:\rk b}
      {\Gamma\vdash((x:T)\to U):\rk{\max(a+1,b)}}.
\]
The other typing rules are the standard rules for variables, abstraction,
application, and $\beta$-conversion. Dependent function types are written
$(x:T)\to U$, and arrows associate to the right. There is no cumulativity,
$\eta$-conversion, or additional proof constant. Omitted lambda annotations
are determined by the displayed types.

Put
\[
 \bot:=((Z:\rk0)\to Z):\rk1,\qquad
 \top:=\rk0\to\rk0:\rk1,\qquad
 \mathsf{id}:=\lambda(Z:\rk0).Z:\top.
\]
Use the fixed-answer proof types
\begin{gather*}
 \nc P:=P\to\bot,\qquad
 P\ac Q:=(P\to Q\to\bot)\to\bot,\\
 \bigvee_c(P_0,\ldots,P_k):=
 (P_0\to\bot)\to\cdots\to(P_k\to\bot)\to\bot,\\
 P\iffc Q:=(P\to Q)\ac(Q\to P).
\end{gather*}
A type $P$ is \emph{stable} when $\nc\nc P\to P$ is inhabited.
Negative types, dependent functions with stable codomain, and these encoded
conjunctions and disjunctions are stable. They support the classical logical
rules on stable proof types; case elimination is into stable targets.
In particular, $P\oc\nc P$ has the proof
$\lambda(k:P\to\bot)(l:(P\to\bot)\to\bot).\,l\,k$.
We never eliminate $\bot$ into an arbitrary positive type.

For a carrier $V$, a domain predicate $\Delta$, and fixed stable atomic
interpretations, use implication as an arrow and translate quantifiers by
\begin{align*}
 (\forall x\,\varphi)^*&:=(x:V)\to\Delta(x)\to\varphi^*(x),\\
 (\exists x\,\varphi)^*&:=
 ((x:V)\to\Delta(x)\to\varphi^*(x)\to\bot)\to\bot.
\end{align*}
All translated formulas are stable. Once domain closure and atomic
congruence are proved, the logical rules and equality substitution extend
structurally to formulas on admitted arguments. All interpretations below
use this one logical encoding. For the relational interpretation,
Section~\ref{sec:relational} establishes completeness of the ground
proof/refutation construction and failure of faithfulness on arbitrary
first-order sentences.

For exact-rank motives, let $J_1=\top$, $J_{k+1}=J_k\to J_k$, with
$j_1=\mathsf{id}$ and $j_{k+1}=\lambda(x:J_k).x$.
If $P:\rk a$, $a\le k$, and $k\ge2$, then
\[
 L_k(P):=J_{k-1}\to P:\rk k.
\]
Abstraction and evaluation at $j_{k-1}$ give maps between $P$ and $L_k(P)$.

\section{Unary domain and equality}\label{sec:unary}
Use the Church numeral carrier
\[
 D:=\bigl((X:\rk1)\to X\to(X\to X)\to X\bigr):\rk3,
 \qquad
 \z:=\lambda X\,z\,s.z,\qquad
 \suc d:=\lambda X\,z\,s.s(d\,X\,z\,s).
\]
Write $\nu_n=\suc^n\z$. Define the two auxiliary, same-type Church operations
\begin{align*}
 \Add\,d\,e&:=\lambda X\,z\,s.\ d\,X\,(e\,X\,z\,s)\,s,\\
 \Mul\,d\,e&:=\lambda X\,z\,s.\ d\,X\,z\,(\lambda y:X.e\,X\,y\,s).
\end{align*}
Both have type $D\to D\to D$. On canonical numerals they compute addition
and multiplication by beta reduction. No eta equation on arbitrary $d:D$
is used.

The domain predicate is the following constructor-closure certificate:
\[
 \Good(d):=\bigl((Q:D\to\rk{10})\to Q(\z)
 \to((u:D)\to Q(u)\to Q(\suc u))\to Q(d)\bigr):\rk{12}.
\]
Zero and successor preserve it by abstraction. For example,
\[
 \mathsf{good}_0:=\lambda Q\,q_0\,q_s.\,q_0,\qquad
 \mathsf{good}_s\,d\,g:=\lambda Q\,q_0\,q_s.\,
 q_s\,d\,(g\,Q\,q_0\,q_s).
\]
Thus the interpreted domain is nonempty and closed under the source
functions. The auxiliary $\Add,\Mul$ occur only inside predicates, so no
closure of $\Good$ under them is required.

Define equality by negative tests:
\[
 \boxed{E(d,e):=\bigl((P:D\to\rk1)\to
       ((\nc P(d))\iffc(\nc P(e)))\bigr):\rk5.}
\]
Its reflexivity, symmetry, and transitivity follow pointwise from the
corresponding rules for logical equivalence. Every $F:D\to D$ respects $E$:
to obtain $E(Fd,Fe)$ from $h:E(d,e)$, specialize $h$ to
$u\mapsto P(Fu)$. This applies to successor and to each argument of the
auxiliary arithmetic operations. Equality substitution for $E$ itself follows
from symmetry and transitivity. Define the arithmetic atoms by
\[
 A^*(d,e,f):=E(\Add\,d\,e,f),\qquad
 M^*(d,e,f):=E(\Mul\,d\,e,f).
\]
They respect $E$ in all arguments. Their ground-input defining schemes are
immediate from the beta equations, including an arbitrary admitted third
argument.

\section{Order from successor-closed unary tests}\label{sec:order}
A certified unary algebra is the fixed telescope
\[
 \begin{aligned}
 \mathcal U={}&(X:\rk1)\,(z:X)\,(s:X\to X)\,(\delta:X\to\rk4)\\
             &(d_0:\delta z)\,
      (d_s:(u:X)\to\delta u\to\delta(su)).
 \end{aligned}
\]
Writing $\mathcal U\to B$ abbreviates the successive dependent arrows for
these binders; it does not introduce a record type.
Put $\ev_{\mathcal U}(d)=d\,X\,z\,s$ and
$G(d):=(\mathcal U\to\delta(\ev_{\mathcal U}(d))):\rk6$.
For $P:X\to\rk1$, write
\[
 \Up_{\mathcal U}(P):=
 (u:X)\to\delta u\to\nc P(u)\to\nc P(su).
\]
Define
\[
 \boxed{\begin{aligned}
 O(d,e):=\bigl(\mathcal U\to(P:X\to\rk1)\to{}&\Up_{\mathcal U}(P)\to{}\\
 &\nc P(\ev_{\mathcal U}(d))\to\nc P(\ev_{\mathcal U}(e))\bigr):\rk6.
 \end{aligned}}
\]
This needs only $d_0,d_s$, and an upward-closure certificate. Preorder laws
and equality compatibility are \emph{derived}, not included as certificates.

\Needspace{7\baselineskip}
\begin{lemma}\label{basic}
$G(\z)$ and $G(d)\to G(\suc d)$ are inhabited, as is
$\Good(d)\to G(d)$. The relation $O$ is reflexive and transitive on all raw
$D$-objects. Moreover, $E(d,e)\to O(d,e)$ and
$G(d)\to O(d,\suc d)$ are inhabited. Both arguments of $O$ respect $E$.
\end{lemma}
\begin{proof}
The two $G$-closure terms use $d_0,d_s$ and the beta fold equations. To
obtain $G(d)$ from $\Good(d)$, instantiate its predicate at
$u\mapsto L_{10}(G(u))$ and evaluate the result at $j_9$.

Reflexivity and transitivity of $O$ are identity and composition on each
test. Given $E(d,e)$, instantiate its predicate with
$u\mapsto P(\ev_{\mathcal U}(u))$ and take the forward implication.
The step law is the upward-closure certificate applied to
$\ev_{\mathcal U}(d)$ and the supplied $G$-evidence. Finally use
$E\Rightarrow O$, symmetry of $E$, and transitivity of $O$ for compatibility.
\end{proof}
Consequently, for each external $i\le j$, there is a closed proof of
$O(\nu_i,\nu_j)$, built by a finite succession of steps.

\section{Finite separation}\label{sec:separation}
\Needspace{7\baselineskip}
\begin{lemma}\label{sep}
For every external $n$, there is a closed term
$\mathsf{no\_back}_n:O(\nu_{n+1},\nu_n)\to\bot$.
Consequently distinct standard numerals are internally $E$-unequal, and
all reversed ground $O$-comparisons have internal refutations.
\end{lemma}
\begin{proof}
Put $c=n+1$, $q=c+1$, and use the rank-one selectors
\[
 F_q:=\underbrace{\rk1\to\cdots\to\rk1}_{q\ \mathrm{arguments}}
           \to\rk1,\qquad
 e_i:=\lambda Z_0\ldots Z_{q-1}.Z_i.
\]
Choose $X=F_q$, $z=e_0$, and the unary table
\[
 s(u):=\lambda Z_0\ldots Z_{q-1}.\
 u\,Z_{\min(0+1,c)}\cdots Z_{\min(c+1,c)}.
\]
Thus $s(e_i)=_\beta e_{\min(i+1,c)}$. Let
\[
 \delta(u):=\bigl((Q:F_q\to\rk3)\to
 Q(e_0)\to\cdots\to Q(e_c)\to Q(u)\bigr):\rk4.
\]
Selector projection gives $d_0$. To construct $d_s$, open the output
$\delta(su)$ at $Q$ and its branch proofs, and specialize the input
certificate to $v\mapsto Q(sv)$. Its branches are the supplied proofs for
the appropriate successor selectors.

Use the unary test
\[
 P(u):=u\,\underbrace{\top\cdots\top}_{c\ \mathrm{entries}}\,\bot.
\]
To prove $\Up(P)$, eliminate $\delta(u)$ into the rank-three motive
$v\mapsto\nc P(v)\to\nc P(sv)$. At $e_c$ both tests are $\bot$ and the
implication is identity. At $e_i$ with $i<c$, a premise
$\top\to\bot$ gives $\bot$ by application to $\mathsf{id}:\top$; abstraction
gives the negative conclusion. 

Specializing $h:O(\nu_c,\nu_n)$ to these data yields
$(\bot\to\bot)\to(\top\to\bot)$. Apply it to $\mathsf{id}_{\bot}$ and
$\mathsf{id}_{\top}$.

If $m>n$, compose $O(\nu_{n+1},\nu_m)$ with a purported
$O(\nu_m,\nu_n)$ and use $\mathsf{no\_back}_n$.
A purported $E(\nu_m,\nu_n)$ supplies this last order proof by
Lemma~\ref{basic}; symmetry handles $m<n$.
\end{proof}
The selectors are certified only through $\delta$; no parametricity claim
about all raw inhabitants of $F_q$ is needed. These are finite separating
observations, not finite models of the entire source theory.

\section{Bounded cases}\label{sec:cases}
Set
\[
 \Case_n(d):=\bigvee_c\bigl(E(d,\nu_0),\ldots,E(d,\nu_n),O(\nu_{n+1},d)\bigr)
       :\rk8.
\]
\Needspace{7\baselineskip}
\begin{lemma}\label{case}
$(d:D)\to\Good(d)\to\Case_n(d)$ is inhabited for every external $n$.
\end{lemma}
\begin{proof}
Use the rank-ten positive invariant
\[
 \Inv_n(d):=(Z:\rk8)\to(G(d)\to\Case_n(d)\to Z)\to Z.
\]
At zero, supply $G(\z)$ and the $E(\z,\nu_0)$ branch. At successor, open
the target continuation and eliminate the child invariant at that same
$Z$, obtaining $G(d)$ and $\Case_n(d)$. Obtain $G(\suc d)$ independently.
To construct the new negative case proof, all branch continuations return
$\bot$. An $i<n$ branch becomes $i+1$ by successor congruence.
The $i=n$ branch gives $E(\suc d,\nu_{n+1})$, hence
$O(\nu_{n+1},\suc d)$. The large branch is composed with
$O(d,\suc d)$. This proves successor closure.

Instantiate $\Good(d)$ with $\Inv_n$. Project its case component by
choosing $Z=\Case_n(d)$ and continuation $\lambda g\,c.c$.
No negative disjunction is eliminated into positive $G$-evidence.
\end{proof}

\section{The relational interpretation}\label{sec:relational}

\Needspace{8\baselineskip}
\begin{theorem}\label{main}
The carrier $D$, admitted domain $\Good$, closed operations $\z,\suc$,
and relations $E,O,A^*,M^*$ interpret the classical relational theory
$\Rt$ in $\PTS$. Every translated axiom has an effectively constructible
closed proof term. The interpretation is closed under classical first-order
derivations: every $\Rt$-derivation of a sentence effectively yields a closed
$\PTS$ proof term of its translation. True ground arithmetic atoms have
proofs, and false ground atoms have internal refutations. The interpretation
is not, however, faithful on arbitrary first-order sentences.
\end{theorem}
\begin{proof}
Use Section~\ref{sec:logic}'s translation with $V=D$ and $\Delta=\Good$.
Domain closure, atomic stability, and congruence have been established;
hence source deductions translate to PTS proofs.

For the first two arithmetic schemes, $\Add\,\nu_m\,\nu_n=_\beta\nu_{m+n}$
and $\Mul\,\nu_m\,\nu_n=_\beta\nu_{mn}$, so both sides of each translated
biconditional are beta-convertible. Numeral inequality is Lemma~\ref{sep}.

For bounded order, apply Lemma~\ref{case}. Under $O(d,\nu_n)$ its large
branch contradicts $\mathsf{no\_back}_n$ by transitivity. Conversely, an
$E(d,\nu_i)$ branch with $i\le n$ gives $O(d,\nu_n)$ using
$E\Rightarrow O$ and the ground comparison. For comparability the finite
branches give $O(d,\nu_n)$ and the large branch gives $O(\nu_n,d)$.

The ground proof/refutation claims follow from the arithmetic beta equations,
reflexivity, ground order, and Lemma~\ref{sep}.
\end{proof}
\begin{remark}
The fixed ranks are $D:3$, $E,A^*,M^*:5$, $G,O:6$, $\delta:4$,
$\Case:8$, $\Inv:10$, and $\Good:12$. The proof uses only $d_0,d_s$
and $\Up(P)$ as observation certificates.
\end{remark}

\paragraph{Completeness on ground arithmetic.}
Every source axiom instance has a proof, and every $\Rt$-derivation translates
using the same fixed definitions. For ground arithmetic atoms, the construction
provides proofs of all true instances and internal refutations of all false
instances. This ground completeness is not reflection of arbitrary source
derivability.

\begin{remark}[Failure of faithfulness]\label{rem:relational-nonfaithful}
The interpretation of $\Rt$ is not faithful: the translation of
$\rho:=\forall x\,(x\le x)$ is provable in $\PTS$, although $\Rt\nvdash\rho$.

\emph{An explicit proof of order reflexivity.}
Its translation is
$\rho^*=(d:D)\to\Good(d)\to O(d,d)$. Expanding every binder in the
observation telescope gives the following closed term:
\[
\begin{aligned}
 \mathsf{orderRefl}:={}&\lambda(d:D)(g:\Good(d))\,\\
 &(X:\rk1)(z:X)(s:X\to X)(\delta:X\to\rk4)\,\\
 &(d_0:\delta z)(d_s:(u:X)\to\delta u\to\delta(su))\,\\
 &(P:X\to\rk1)\,\\
 &(p_{\uparrow}:(u:X)\to\delta u\to
       (P(u)\to\bot)\to(P(su)\to\bot))\,\\
 &(h:P(d\,X\,z\,s)\to\bot).\;h.
\end{aligned}
\]
Thus
\begin{equation}
 \mathsf{orderRefl}:(d:D)\to\Good(d)\to O(d,d).
 \label{eq:order-reflexivity}
\end{equation}
Indeed, after the binders are opened, the remaining goal is
$P(d\,X\,z\,s)\to\bot$, exactly the type of $h$. The term simply returns
$h$; it uses neither the domain proof $g$ nor the three observation
certificates $d_0,d_s,p_{\uparrow}$. Removing the binder $g$ also gives
$(d:D)\to O(d,d)$, so the proof works on every raw presentation.

\emph{A relational countermodel.}
Take $N^+=\mathbb N\sqcup\{\omega\}$, where $\omega$ is one new element.
Interpret $0$ and successor as usual on $\mathbb N$, put $S\omega=\omega$,
and use actual equality. For $x,y\in\mathbb N$, let $A(x,y,z)$ hold exactly
when $z=x+y$, and let $M(x,y,z)$ hold exactly when $z=xy$. When at least one
of $x,y$ is $\omega$, let each predicate hold exactly when $z=0$. Define
\[
 x\le^+ y\quad\Longleftrightarrow\quad
 \bigl(x,y\in\mathbb N\ \text{and}\ x\le_{\mathbb N}y\bigr)
 \ \text{or}\ \bigl(x\in\mathbb N\ \text{and}\ y=\omega\bigr).
\]
The ground graph schemes hold for every third argument, and distinct numerals
are unequal. For every standard $n$, the downset of $n$ is exactly
$\{0,\ldots,n\}$, proving both directions of \eqref{rtilde:bound}.
Every standard number is comparable with $n$, and $n\le^+\omega$, proving
\eqref{rtilde:compare}. Hence this is a model of $\Rt$.

However, $\omega\not\le^+\omega$, so $\rho$ is false. By soundness of
classical first-order logic,
\[
 \Rt\nvdash\rho,
 \qquad\text{whereas}\qquad \PTS\vdash\mathsf{orderRefl}:\rho^*.
\]
Thus the relational interpretation preserves derivability but does not
reflect it. The countermodel is a model of the source theory, not a
set-theoretic model of $\PTS$. There is no conflict with ground completeness:
$\rho$ is universally quantified, and its failure occurs at a non-numeral
element.
\end{remark}

\section{Relation to the original presentation of \texorpdfstring{$\R$}{R}}\label{sec:nonrelational}
The original Tarski--Mostowski--Robinson arithmetic $\R$~\cite{tmr} uses
$\{0,S,+,\cdot,\le\}$ and the following schemes, again with universal closure:
\begin{align}
 \bar m+\bar n&=\overline{m+n}, &
 \bar m\cdot\bar n&=\overline{mn},
     \tag{R1--R2}\label{r:arithmetic}\\
 \bar m&\ne\bar n\quad(m\ne n), &&
     \tag{R3}\label{r:neq}\\
 x\le\bar n&\longrightarrow\bigvee_{i=0}^{n}x=\bar i, &&
     \tag{R4}\label{r:bound}\\
 x\le\bar n&\lor\bar n\le x. &&
     \tag{R5}\label{r:compare}
\end{align}
Requiring $A,M$ in $\Rt$ to be total, single-valued graphs makes them
interchangeable with the function symbols $+,\cdot$. This functional
strengthening is definitionally equivalent to $\R$. Even without that
strengthening, the two theories have the same interpretability degree.

\Needspace{8\baselineskip}
\begin{proposition}\label{prop:graphs}
The classical theories $\Rt$ and $\R$ are mutually interpretable.
Consequently $\PTS$ interprets $\R$ by definable arithmetic graphs.
\end{proposition}
\begin{proof}
In $\R$, interpret $A(x,y,z)$ by $x+y=z$ and $M(x,y,z)$ by $xy=z$.
The ground arithmetic equations give the two graph schemes.
For $i<j$, an assumption $\bar j\le\bar i$ contradicts
\eqref{r:bound} and numeral separation; \eqref{r:compare} therefore yields
$\bar i\le\bar j$. The case $i=j$ follows from \eqref{r:compare} directly.
Thus $\R$ proves all true ground comparisons, and equality substitution
supplies the reverse implication in \eqref{rtilde:bound}.
The remaining relational schemes are immediate.

For the other direction define, in $\Rt$,
\begin{align*}
 U_A(x,y,z)&:=A(x,y,z)\land
                \forall w\bigl(A(x,y,w)\to w=z\bigr),\\
 F_A(x,y,z)&:=U_A(x,y,z)\lor
       \bigl(z=0\land\neg\exists v\,U_A(x,y,v)\bigr).
\end{align*}
For fixed $x,y$, either a unique $A$-output exists or it does not. In the
first case $F_A$ selects that output; in the second it selects zero.
Hence $F_A$ is provably total and single-valued. At standard inputs,
\eqref{rtilde:add} forces its value to be $\overline{m+n}$.
Define $F_M$ identically with $M$ in place of $A$; its standard value is
$\overline{mn}$. These graphs interpret $+,\cdot$; keep $0,S,=,\le$
unchanged. All five schemes of $\R$ follow.

Composing this graph interpretation with Theorem~\ref{main} gives the
claimed interpretation in $\PTS$. The composition uses the encoded
classical quantifiers to express values and uniqueness; it does not
assert $\Good$-closure of the raw Church operators $\Add,\Mul$.
\end{proof}
The conclusion is mutual interpretability, not bi-interpretability.
In particular, the totalized graphs need not coincide with the original
relations $A,M$ on arbitrary arguments.


\section{A direct interpretation with arithmetic operators}\label{sec:operators}
The following construction gives the source operations as closed terms
with domain-closure certificates. This section defines its domain and
observations; Section~\ref{sec:arith-finite} proves finite separation and
bounded cases, and Section~\ref{sec:functional} verifies the interpretation.
We use the logic, selectors, and invariant method of
Sections~\ref{sec:logic}--\ref{sec:cases}. Subscript $a$ marks the
arithmetic-expression domain and its observations.

\subsection{The arithmetic-expression domain}
Set
\[
 \Da:=\bigl((X:\rk1)\to X\to(X\to X)\to(X\to X\to X)
                    \to(X\to X\to X)\to X\bigr):\rk3.
\]
For $\xi=(X,z,s,a,m)$ in this telescope, write
$\operatorname{fold}_{\xi}(d):=d\,X\,z\,s\,a\,m$. Define
\begin{align*}
 \za&:=\lambda\xi.\,z,&
 \sa\,d&:=\lambda\xi.\,s(\operatorname{fold}_{\xi}(d)),\\
 \adda\,d\,e&:=\lambda\xi.\,a(\operatorname{fold}_{\xi}(d))
                          (\operatorname{fold}_{\xi}(e)),&
 \ma\,d\,e&:=\lambda\xi.\,m(\operatorname{fold}_{\xi}(d))
                          (\operatorname{fold}_{\xi}(e)).
\end{align*}
Here $\lambda\xi$ abbreviates five successive abstractions, not a record
argument. Let $\nua n=\sa^n\za$. The binary operations build expression
presentations; their arithmetic equations will be proved in $\Ea$, not
asserted as $\beta$-equations on canonical numerals.

For $Q:\Da\to\rk{10}$ abbreviate
\begin{align*}
 \mathsf{Cl}_s(Q)&:=(u:\Da)\to Q(u)\to Q(\sa u),\\
 \mathsf{Cl}_f(Q)&:=(u:\Da)\to(v:\Da)\to Q(u)\to Q(v)\to Q(f\,u\,v)
       \quad(f=\adda,\ma).
\end{align*}
Use
\[
 \Gooda(d):=\bigl((Q:\Da\to\rk{10})\to Q(\za)
 \to\mathsf{Cl}_s(Q)\to\mathsf{Cl}_{\adda}(Q)
 \to\mathsf{Cl}_{\ma}(Q)\to Q(d)\bigr):\rk{12}.
\]
The four constructor certificates are obtained by abstraction. For example,
from $g:\Gooda(d)$ and $h:\Gooda(e)$, addition closure is
\[
 \lambda Q\,q_0\,q_s\,q_a\,q_m.\;
 q_a\,d\,e\,(g\,Q\,q_0\,q_s\,q_a\,q_m)
             (h\,Q\,q_0\,q_s\,q_a\,q_m).
\]
Thus the admitted domain is nonempty and closed under all four source
operations, without assuming induction on arbitrary raw $\Da$-objects.

\subsection{Certified arithmetic observations}\label{sec:arith-obs}
An arithmetic observation has data
\[
 X:\rk1,\quad z:X,\quad s:X\to X,\quad a,m:X\to X\to X,
 \quad\delta:X\to\rk4,\quad R_0,C_0:X\to X\to\rk1.
\]
Write $u\simeq v:=R_0(u,v)\to\bot$ and
$u\preceq v:=C_0(u,v)\to\bot$. Its remaining telescope parameters
certify the following laws. Each object variable is immediately guarded
by $\delta$; for example $H(u,v)$ abbreviates
$(u:X)\to\delta(u)\to(v:X)\to\delta(v)\to H(u,v)$.
Each law is a separate proof argument.

\emph{Domain and congruence.} The predicate $\delta$ contains $z$ and is
closed under $s,a,m$. The relation $\simeq$ is an equivalence respected by
all three operations. The relation $\preceq$ is reflexive and transitive,
and respects $\simeq$ in both arguments:
\[
 u\simeq u'\to v\simeq v'\to u\preceq v\to u'\preceq v'.
\]
\emph{Arithmetic equations.}
\begin{align}
 a(z,u)&\simeq u,& a(su,v)&\simeq s(a(u,v)),\notag\\
 m(z,u)&\simeq z,& m(su,v)&\simeq a(v,m(u,v)),\qquad m(u,z)\simeq z.
 \label{eq:arith-laws}
\end{align}
\emph{Order propagation}, with $1_X=sz$:
\begin{gather}
 u\preceq su,\qquad u\preceq a(u,v),\qquad v\preceq a(u,v),\notag\\
 1_X\preceq v\to u\preceq m(u,v),\qquad
 1_X\preceq u\to v\preceq m(u,v).
 \label{eq:arith-order}
\end{gather}
There are 23 individual certificates, each of rank at most five. They
are bound parameters, not axioms of $\PTS$.

Denote the complete telescope by $\mathcal A$ and put
$\ev_{\mathcal A}(d)=d\,X\,z\,s\,a\,m$. Define
\begin{align*}
 \Ea(d,e)&:=\mathcal A\to
          \ev_{\mathcal A}(d)\simeq\ev_{\mathcal A}(e),\\
 \Oa(d,e)&:=\mathcal A\to
          \ev_{\mathcal A}(d)\preceq\ev_{\mathcal A}(e),\\
 \Ga(d)&:=\mathcal A\to\delta(\ev_{\mathcal A}(d)).
\end{align*}
All three have rank six; $\Ea,\Oa$ are stable and $\Ga$ is positive.
The fold equations and the domain certificates make $\Ga$ constructor
closed. Instantiating $\Gooda(d)$ at $u\mapsto L_{10}(\Ga(u))$ gives
$\Gooda(d)\to\Ga(d)$. Denote the resulting closed map by
$\gamma_a:(d:\Da)\to\Gooda(d)\to\Ga(d)$.

Opening the observation telescope transfers all its guarded laws to
$\Ga$-objects. In particular, $\Ea$ is an equivalence respected by
$\sa,\adda,\ma$ and $\Oa$ is a congruent preorder. External induction on
standard numeral indices gives closed proofs of
\begin{equation}
 \Ea(\adda\,\nua m\,\nua n,\nua{m+n}),\quad
 \Ea(\ma\,\nua m\,\nua n,\nua{mn}),\quad
 \Oa(\nua i,\nua j)\quad(i\le j).
 \label{eq:arith-ground}
\end{equation}
For the first two use \eqref{eq:arith-laws}; for the last use reflexivity,
the successor step, and transitivity. No internal arithmetic induction
is assumed.

\section{Finite separation and bounded cases}\label{sec:arith-finite}
For external $K\ge0$, set $c=K+1$. On $\{0,\ldots,c\}$ use
\[
 z=0,\qquad s(i)=\min(i+1,c),\qquad
 a(i,j)=\min(i+j,c),\qquad m(i,j)=\min(ij,c),
\]
with ordinary equality and order. These tables satisfy the observation
laws. In the multiplication recursion, $j=0$ makes both sides zero; if
$j>0$, overflow makes both sides equal $c$. For propagation,
$j\ge1$ and $i\le c$ imply $\min(ij,c)\ge i$. Addition and the remaining
laws are checked directly. In particular, $m(c,0)=0$.

Implement the tables using $q=K+2$ and the selectors $F_q,e_i$ of
Section~\ref{sec:separation}. A unary table is implemented as there; a
binary table $b$ is
\[
 \widehat b(u,v):=\lambda Z_0\ldots Z_{q-1}.\;
 u\,(v\,Z_{b(0,0)}\cdots Z_{b(0,q-1)})\cdots
   (v\,Z_{b(q-1,0)}\cdots Z_{b(q-1,q-1)}).
\]
For a relation table, omit the outer abstractions and replace each
selected $Z_{b(i,j)}$ by $\bot$ at a true entry and $\top$ at a false
entry. This gives $R_0,C_0:X\to X\to\rk1$; their negations are the
required proof relations.

Use the same $\delta_q$ as in Section~\ref{sec:separation}.
Selector introduction certifies $z$. To certify an operation's closure,
open the output $\delta_q$ at a rank-three motive $Q$ and eliminate the
input certificates into $Q$ composed with the operation. For relation
laws, open all domain guards and eliminate into the remaining goal,
whose rank is at most three; apply $L_3$ when its rank is smaller.
The resulting branches are precisely the finite table checks. A false
relational premise yields $\bot$ on application to $\mathsf{id}:\top$;
negative conclusions follow by abstraction. Thus every finite table
above has a closed implementation of the \emph{entire} telescope
$\mathcal A$, including its 23 proof arguments.

At cutoff $K=n$, the folds of $\nua{n+1}$ and $\nua n$ are in strictly
reversed order. Specializing a proposed $\Oa$-proof gives $\top\to\bot$;
applying it to $\mathsf{id}$ constructs
\[
 \mathsf{no\_back}^{a}_n:
 \Oa(\nua{n+1},\nua n)\to\bot.
\]
Ground comparisons and congruence then give all reversed-order and
unequal-numeral refutations exactly as in Lemma~\ref{sep}.

\Needspace{12\baselineskip}
\paragraph{Bounded cases.}
Define
\[
 \Casea_n(d):=\bigvee_c\bigl(
 \Ea(d,\nua0),\ldots,\Ea(d,\nua n),\Oa(\nua{n+1},d)\bigr):\rk8.
\]
\Needspace{7\baselineskip}
\begin{lemma}\label{lem:arith-cases}
For every external $n$, $(d:\Da)\to\Gooda(d)\to\Casea_n(d)$ is inhabited.
\end{lemma}
\begin{proof}
Use the same positive invariant template as in Lemma~\ref{case}:
\[
 \Inva_n(d):=(Z:\rk8)\to(\Ga(d)\to\Casea_n(d)\to Z)\to Z.
\]
It has rank ten. Each constructor combines the children's positive
$\Ga$-evidence independently of their cases. Open the parent's case
continuations before splitting the child cases, so every branch returns
$\bot$. The four closure arguments are as follows.

\emph{Zero and successor.} Reflexivity gives the zero case. A successor
of an $i<n$ case has index $i+1$; at $i=n$ its equality to $\nua{n+1}$
gives the large branch. A large child remains large by the successor
step. This is the argument of Lemma~\ref{case} using the transferred laws.

\emph{Addition.} Finite cases $i,j$ give equality to $\nua{i+j}$ by
congruence and \eqref{eq:arith-ground}. If $i+j\le n$, take that finite
branch; otherwise ground order gives the large branch. A large child
bounds the sum below by \eqref{eq:arith-order}.

\emph{Multiplication.} Two finite cases similarly give equality to
$\nua{ij}$. A zero factor makes the product equal to zero, using the left
or right zero equation as appropriate. A nonzero finite factor is at
least one, and a large factor is also at least one since $1\le n+1$.
Thus a large factor paired with a nonzero factor bounds the product
below by the multiplication propagation laws. This covers both mixed
cases and the case of two large factors.

Supply the resulting positive $\Ga$-proof and negative case proof to the
parent continuation. These are the four closure proofs of $\Inva_n$.
Instantiate $\Gooda(d)$ at that motive and project the case component by
taking $Z=\Casea_n(d)$. No negative case proof is used to manufacture
positive domain evidence.
\end{proof}

\section{The functional interpretation}\label{sec:functional}

\Needspace{8\baselineskip}
\begin{theorem}\label{thm:direct-R}
The carrier $\Da$, admitted domain $\Gooda$, closed operations
$\za,\sa,\adda,\ma$, and relations $\Ea,\Oa$ directly interpret classical
$\R$ in $\PTS$. Every translated axiom has an effectively constructible
closed proof term. The interpretation is closed under classical first-order
derivations: every $\R$-derivation of a sentence effectively yields a closed
$\PTS$ proof term of its translation. True ground arithmetic atoms have
proofs, and false ground atoms have internal refutations. The interpretation
is not, however, faithful on arbitrary first-order sentences.
\end{theorem}
\begin{proof}
Translate terms compositionally using these four operations, and use
Section~\ref{sec:logic}'s logical translation with $V=\Da$ and
$\Delta=\Gooda$. Domain closure, stability, and congruence were proved in
Section~\ref{sec:arith-obs}; therefore the logical and equality rules
preserve typing of translated proofs.

The equations \eqref{r:arithmetic} are \eqref{eq:arith-ground}, and
\eqref{r:neq} follows from finite separation. For \eqref{r:bound}, assume
$g:\Gooda(d)$ and $h:\Oa(d,\nua n)$ and apply
Lemma~\ref{lem:arith-cases}. A finite branch gives the required equality;
the large branch yields $\Oa(\nua{n+1},\nua n)$ by transitivity and is
refuted by $\mathsf{no\_back}^{a}_n$. For \eqref{r:compare}, a finite
branch gives $\Oa(d,\nua n)$, and the large branch gives
$\Oa(\nua n,d)$ by composing with the ground successor comparison.
These eliminations have stable targets.

Finally, external evaluation of a closed arithmetic term $t$ with value
$k$ constructs $\Ea(t^\dagger,\nua k)$ by induction on $t$, using
congruence and \eqref{eq:arith-ground}. Ground comparisons and finite
separation provide the required proofs or refutations of every closed
equation or comparison. Every external induction builds a finite PTS
term; the interpretation itself is independent of the axiom indices.
\end{proof}

\Needspace{14\baselineskip}
\paragraph{Order reflexivity in the direct interpretation.}
Name the order-reflexivity certificate in $\mathcal A$ by
\[
 r_{\preceq}:(u:X)\to\delta(u)\to(C_0(u,u)\to\bot).
\]
The map $\gamma_a$ of Section~\ref{sec:arith-obs} gives the closed term
\[
 \mathsf{orderRefl}_a:=
 \lambda(d:\Da)(g:\Gooda(d))\,\mathcal A.\;
 r_{\preceq}\,\ev_{\mathcal A}(d)\,
               (\gamma_a\,d\,g\,\mathcal A).
\]
Here $\lambda\mathcal A$ expands all data and certificate binders,
including $r_{\preceq}$. The last argument proves
$\delta(\ev_{\mathcal A}(d))$, enabling the guarded reflexivity law. Hence
\[
 \mathsf{orderRefl}_a:\rho^\dagger,
 \qquad \rho^\dagger:=(d:\Da)\to\Gooda(d)\to\Oa(d,d).
\]
The relational countermodel at the end of Section~\ref{sec:relational}
expands to a model of $\R$: take $+$ and $\cdot$ to be the total functions
whose graphs are its predicates $A,M$. They have the required values on
standard inputs. Numeral separation and the two order schemes are unchanged,
but $\omega\not\le^+\omega$ still holds. Thus $\R\nvdash\rho$, and
$\mathsf{orderRefl}_a$ shows that the direct interpretation is not faithful.
The same applies to the graph-based interpretation from
Proposition~\ref{prop:graphs}: it keeps the unary domain and $O$ unchanged,
so \eqref{eq:order-reflexivity} supplies its proof of $\rho^*$.

\paragraph{Scope and nontriviality.}
The ground proof/refutation results remain valid. External uninhabitation of
$\bot$ follows from strong normalization, subject reduction, and exclusion of
normal terms of type $X$ in the context $X:\rk0$. No set model of $\PTS$
is used. These constructions are not a completed Rocq formalization.

\begin{thebibliography}{1}\small
\bibitem{tmr}
A. Tarski, A. Mostowski, and R. M. Robinson.
\emph{Undecidable Theories}. North-Holland, 1953.
\end{thebibliography}
\end{document}
